\paragraph{Corrections to ``Ensemble singulier et supersingulier'' (cote 25, sites 165--167).}
The faint opening on site 165 and the clearer page headed ``EGA IV par ?'' on site 166 introduce a further set of corrections to the earlier notes. Site 165 retains extensive local doubts; its pale impression is not silently replaced by the clearer wording on site 166. The opening instructions on site 166 change $S$ to $Y$ on manuscript pages 14 and 15 and object to the ``i.e.'' on page 15, line 10. The subsequent argument addresses the generalization of proposition 7 mentioned on pages 16 and 17. Its diagram combines the conormal sheaf of $V^1$ in $X$ with the relative differential sheaves for $X/Y$, $X/S$, $Y/S$ and $V^1/S$. It explicitly invokes the lemma on manuscript page 7 and compares the zero scheme $V(a)$ with $V(D)$. The nested immersions $V^1\subset V^o\subset X$ provide the displayed conormal sequence, for which the text leaves a reference to be supplied. The two following identifications use $E$ and $\mathrm{Hom}(W,K)$; $W$ and $K$ are described by the cokernel and kernel of the displayed map $v^1$. Site 167 refers back to manuscript page 5 for the first identification and to the attached differential complements to ``variation du rang d'un homomorphisme de Modules loc libres'' for the second. It then separates two checks: identifying $\mu i$ with $v^1$, and identifying the induced $\mu'$ with the homomorphism of the quadratic form from the bottom of page 15. The author leaves these checks to Dieudonné and expressly raises a possible sign. The guide therefore records the route of the proposed verification and its dependencies, without supplying that sign or claiming that the omitted checks were performed. The written comparison signs, the occurrences of $V(u)$, and the unsettled overwritten cells remain in the text. See \href{https://grothendieck.umontpellier.fr/web/viewer.html?file=../25.pdf#page=166}{site 165}, \href{https://grothendieck.umontpellier.fr/web/viewer.html?file=../25.pdf#page=167}{site 166} and \href{https://grothendieck.umontpellier.fr/web/viewer.html?file=../25.pdf#page=168}{site 167}.

\paragraph{The proposed intrinsic descriptions and their unfinished verification (sites 168--171).}
Site 168 revisits the change of viewpoint mentioned in the author's last letter. It starts with a closed subscheme $V^o=V(J)$ of $X$, the conormal Module $N=J/J^2$, and a local presentation through a finite locally free Module $E$. The discussion asks for the independence of $V^1$, $V''$ and $V^2$ from that presentation, while retaining the dependence on the written ambient data and on $m$. It describes $V^1$ and $V''$ by the exterior-power maps labelled $(*_m)$ and $(*_{m-1})$. The opening formula on site 169 has an unmatched parenthesis in the scan, which has been retained. For $V^2$, site 169 expressly says that the invariance has not really been verified and gives it conjectural status. The proposed route compares the maps from $E$ and from $N$, names their common cokernel $W$, and introduces the kernel $k$ of the intrinsic map as a quotient of $K$. It asks for the previously defined $M:K\to\mathrm{Sym}^2(W)$ to factor through a presentation-independent $m:k\to\mathrm{Sym}^2(W)$. The rank phrase beside $W$ contains a disputed cancellation and is left unresolved. Site 170 continues this proposed route with a representing functor, whose slashed-loop symbol remains locally ambiguous, then describes the intended complementary locus in terms of flatness, geometric regularity or ordinary quadratic singularities, and transverse regularity. The displayed prose also has an actual blank between $E$ and $J$; an arrow has not been supplied. Site 171 discusses the inverse-image situation, the available global surjection from $E^o$ to the conormal Module, and the resulting discriminant expression. Its handwritten insertion before ``à priori'' reads ``ici''. The final paragraph proposes adding the invariance discussion in small print after the locally free Module case. It weighs paragraph 19 against paragraph 22, explicitly locating the differential argument used on manuscript pages 18 and 19 with the Grassmannian material, and suggests a matrix calculation as an alternative exposition. These are the author's stated proposals and reservations. See \href{https://grothendieck.umontpellier.fr/web/viewer.html?file=../25.pdf#page=169}{site 168}, \href{https://grothendieck.umontpellier.fr/web/viewer.html?file=../25.pdf#page=170}{site 169}, \href{https://grothendieck.umontpellier.fr/web/viewer.html?file=../25.pdf#page=171}{site 170} and \href{https://grothendieck.umontpellier.fr/web/viewer.html?file=../25.pdf#page=172}{site 171}.
